\chapter{Time Evolution, Conservation, and Coupled Spins}

A quantum state is not a fixed point in Hilbert space. As time passes it moves, yet it must move without losing normalization, probabilities, or the distinguishability encoded in inner products. Those requirements nearly determine the form of quantum dynamics. Linearity gives a time-evolution operator, probability conservation makes that operator unitary, and an infinitesimal unitary transformation is generated by a Hermitian operator. That operator is the Hamiltonian. We shall derive the Schr\"odinger equation from these statements, see how energy eigenstates accumulate phase, and then turn the formalism into two concrete calculations: precession of one spin in a magnetic field and oscillation between two states of a coupled spin pair.

\section{Linear Motion in Hilbert Space}

Let the state at time \(t\) be \(\ket{\psi(t)}\). We assume that its relation to the state at time zero is linear:
\begin{equation}
\ket{\psi(t)}=U(t)\ket{\psi(0)}.
\end{equation}
The linearity is a postulate, not a theorem derived from something more elementary. It is an experimentally successful part of quantum mechanics, and attempts to replace it generally endanger either probability or conservation laws. In a chosen basis, \(U(t)\) is simply a time-dependent matrix.

The dual vector obeys the conjugate rule. If
\begin{equation}
\ket{\phi(t)}=U(t)\ket{\phi(0)},
\end{equation}
then
\begin{equation}
\bra{\phi(t)}=\bra{\phi(0)}U^\dagger(t).
\end{equation}
For a matrix, the dagger means complex conjugation followed by transposition. The bra therefore carries the same evolution law as the ket, expressed in the dual space.

\subsection{Probability Conservation Implies Unitarity}

A normalized state represents total probability one:
\begin{equation}
\braket{\psi(t)}{\psi(t)}=1.
\end{equation}
Time evolution must preserve this norm. In fact, a linear map that preserves the norm of every vector also preserves every inner product. Geometrically, preserving the lengths of two vectors and their difference preserves the triangle formed by them and hence their angle. In a complex vector space the same conclusion follows by also considering combinations with a factor of \(i\).

Thus, for arbitrary \(\ket{\phi}\) and \(\ket{\psi}\),
\begin{equation}
\braket{\phi(t)}{\psi(t)}
=
\braket{\phi(0)}{\psi(0)}.
\end{equation}
Substituting the evolution laws gives
\begin{equation}
\bra{\phi(0)}U^\dagger(t)U(t)\ket{\psi(0)}
=
\braket{\phi(0)}{\psi(0)}.
\end{equation}
Since this equation holds for every pair of vectors,
\begin{equation}
\boxed{U^\dagger(t)U(t)=I}.
\end{equation}
Such an operator is unitary. The scalar analogue is a complex number of unit magnitude,
\begin{equation}
e^{i\theta}\bigl(e^{i\theta}\bigr)^*=1.
\end{equation}
A unitary operator is the matrix generalization of a pure phase: it can rotate a state through Hilbert space, but it cannot stretch or shrink it.

\section{The Infinitesimal Generator}

Finite evolution is built from very small time steps. Let the elapsed time be \(\epsilon\), small enough that terms of order \(\epsilon^2\) may be ignored. At zero elapsed time nothing happens, so
\begin{equation}
U(0)=I.
\end{equation}
The leading correction must be proportional to \(\epsilon\). We write it in the conventional form
\begin{equation}
\boxed{U(\epsilon)=I-\frac{i\epsilon}{\hbar}H}.
\end{equation}
Only the factor of \(\epsilon\) is forced at this stage. The \(i\), \(\hbar\), and the name \(H\) are chosen so that the object left at the end has the standard normalization and is Hermitian.

\begin{classroomqa}
\lecturetimestamp{00:12:46}
\audiencequestion{Why is Planck's constant such a remarkably small number in familiar units?}
\lecturerresponse{The numerical size reflects the units chosen by human beings. Chemists selected macroscopic amounts that fit in hand-sized vessels, while atoms are vastly smaller than people. In units adapted to those macroscopic scales, an atomic quantum of action must look tiny. This is not a separate dynamical mystery; it is largely the historical and human scale built into the units.}
\end{classroomqa}

Taking the Hermitian conjugate,
\begin{equation}
U^\dagger(\epsilon)
=I+\frac{i\epsilon}{\hbar}H^\dagger.
\end{equation}
Now multiply and retain only first-order terms:
\begin{align}
U^\dagger(\epsilon)U(\epsilon)
&=
\left(I+\frac{i\epsilon}{\hbar}H^\dagger\right)
\left(I-\frac{i\epsilon}{\hbar}H\right) \\
&=
I+\frac{i\epsilon}{\hbar}\left(H^\dagger-H\right)
+O(\epsilon^2).
\end{align}
Unitarity requires this product to equal \(I\), and therefore
\begin{equation}
\boxed{H^\dagger=H}.
\end{equation}
The infinitesimal generator of time evolution is Hermitian. Since Hermitian operators represent observables, every dynamical system possesses an observable that generates its motion in time. This is the Hamiltonian. We use \(H\) for the operator and reserve \(E\) for its possible measured values, the eigenvalues.

\begin{classroomqa}
\lecturetimestamp{00:18:34}
\audiencequestion{Why should this generator be called energy? Is energy not \(p^2/2m\), or the capacity to do work?}
\lecturerresponse{Those are useful descriptions in special classical settings, not a general definition. The expression \(p^2/2m\) is only the kinetic energy of a nonrelativistic free particle, and it does not define such things as energy stored in an electromagnetic field. For an abstract system the first universal feature to demand of energy is conservation. We shall shortly prove that the expectation value of \(H\) is conserved, and its eigenstate phases will also satisfy \(E=\hbar\omega\). Those two facts identify the generator with energy.}
\end{classroomqa}

\section{From a Small Step to the Schr\"odinger Equation}

The infinitesimal update is already a differential equation in disguise. At an arbitrary origin of time,
\begin{equation}
\frac{d}{dt}\ket{\psi(t)}
=
\lim_{\epsilon\to0}
\frac{\ket{\psi(t+\epsilon)}-\ket{\psi(t)}}{\epsilon}.
\end{equation}
The board first writes this at \(t=0\); shifting the origin gives the general form above.

\begin{figure}[htbp]
\centering
\includegraphics[width=0.72\textwidth]{lecture_09_figure_03.png}
\caption{The difference quotient that turns infinitesimal evolution into a differential equation.}
\lectureframe{00:22:15}
\end{figure}

Use
\begin{equation}
\ket{\psi(t+\epsilon)}
=
\left(I-\frac{i\epsilon}{\hbar}H\right)\ket{\psi(t)}.
\end{equation}
Then
\begin{align}
\frac{d}{dt}\ket{\psi(t)}
&=
\lim_{\epsilon\to0}
\frac{\left(I-\frac{i\epsilon}{\hbar}H\right)\ket{\psi(t)}
-\ket{\psi(t)}}{\epsilon} \\
&=
-\frac{i}{\hbar}H\ket{\psi(t)}.
\end{align}
Equivalently,
\begin{equation}
\boxed{
i\hbar\frac{d}{dt}\ket{\psi(t)}
=H\ket{\psi(t)}
}.
\end{equation}
This is the Schr\"odinger equation. Once \(H\) and the state at one instant are known, it determines the state at later times. Probabilities for later measurements then follow from the ordinary measurement rules.

\section{Energy Eigenstates and Rotating Phase}

The simplest solution begins with an eigenvector of the Hamiltonian:
\begin{equation}
H\ket{\psi_E}=E\ket{\psi_E}.
\end{equation}

\begin{figure}[htbp]
\centering
\includegraphics[width=0.70\textwidth]{lecture_09_figure_04.png}
\caption{The energy-eigenvalue equation on the blackboard.}
\lectureframe{00:28:22}
\end{figure}

Since \(H\) sends \(\ket{\psi_E}\) back into the same one-dimensional direction, time evolution cannot mix it with another eigenvector. It can only multiply it by a scalar:
\begin{equation}
\ket{\psi_E(t)}=f(t)\ket{\psi_E}.
\end{equation}
Substitution into the Schr\"odinger equation gives
\begin{align}
i\hbar\,\dot f(t)\ket{\psi_E}
&=
f(t)H\ket{\psi_E} \\
&=
E f(t)\ket{\psi_E}.
\end{align}
Therefore
\begin{equation}
\dot f(t)=-\frac{iE}{\hbar}f(t),
\qquad
f(t)=e^{-iEt/\hbar},
\end{equation}
and
\begin{equation}
\boxed{
\ket{\psi_E(t)}
=e^{-iEt/\hbar}\ket{\psi_E(0)}
}.
\end{equation}
An energy eigenstate is stationary in the sense that its ray in Hilbert space does not move. Its ket nevertheless accumulates phase.

The exponential may be written
\begin{equation}
e^{-iEt/\hbar}
=
\cos\!\left(\frac{Et}{\hbar}\right)
-i\sin\!\left(\frac{Et}{\hbar}\right).
\end{equation}
The minus sign and the absence of an extra \(i\) inside the sine were corrected during the blackboard calculation. The angular frequency is
\begin{equation}
\omega=\frac{E}{\hbar},
\qquad
\boxed{E=\hbar\omega}.
\end{equation}
Equivalently, using ordinary frequency \(\nu=\omega/(2\pi)\), one has \(E=h\nu\).

\begin{figure}[htbp]
\centering
\includegraphics[width=0.78\textwidth]{lecture_09_figure_05.png}
\caption{The phase solution and the energy--frequency relation are developed on adjacent boards.}
\lectureframe{00:36:19}
\end{figure}

\begin{classroomqa}
\lecturetimestamp{00:34:43}
\audiencequestion{Why were the factors \(i\) and \(\hbar\) inserted in the infinitesimal operator before the calculation required them?}
\lecturerresponse{They are conventions chosen with the answer in view. Without the \(i\), unitarity would make the generator anti-Hermitian. Multiplication by \(i\) turns it into a Hermitian observable with real eigenvalues. The factor \(\hbar\) gives those eigenvalues the units and normalization convention of energy. The physical statement is unitarity; these factors put its generator into standard form.}
\end{classroomqa}

\begin{classroomqa}
\lecturetimestamp{00:36:18}
\audiencequestion{Does \(E=\hbar\omega\) belong only to electromagnetic waves and photons?}
\lecturerresponse{No. Photons are one important case, but the relation is more general. Any Hamiltonian eigenstate acquires a phase with angular frequency \(E/\hbar\). Connecting the frequency of an electromagnetic field mode to the frequency of its photon requires more physics, but the energy--phase relation itself already follows from quantum time evolution.}
\end{classroomqa}

\section{Superposition in the Energy Basis}

A single energy eigenstate changes only by an overall phase, and an overall phase cannot affect a probability. Interesting dynamics begins when a state contains several energy eigenvectors. Since a Hermitian Hamiltonian has a complete orthonormal eigenbasis,
\begin{equation}
H\ket{\psi_n}=E_n\ket{\psi_n},
\qquad
\ket{\psi(0)}=\sum_n\alpha_n\ket{\psi_n}.
\end{equation}

\begin{figure}[htbp]
\centering
\includegraphics[width=0.72\textwidth]{lecture_09_figure_06.png}
\caption{Expansion of an arbitrary state in Hamiltonian eigenvectors.}
\lectureframe{00:38:35}
\end{figure}

Linearity permits each component to evolve separately:
\begin{equation}
\boxed{
\ket{\psi(t)}
=
\sum_n\alpha_n e^{-iE_n t/\hbar}\ket{\psi_n}
}.
\end{equation}
High-energy components wind rapidly around the complex plane; low-energy components wind more slowly. What changes measurable probabilities is not a phase shared by the whole state, but the changing relative phases between components with different \(E_n\).

In the energy basis,
\begin{equation}
\begin{aligned}
H&=
\begin{pmatrix}
E_1&0&\cdots\\
0&E_2&\cdots\\
\vdots&\vdots&\ddots
\end{pmatrix},\\[0.6em]
U(t)&=
\begin{pmatrix}
e^{-iE_1t/\hbar}&0&\cdots\\
0&e^{-iE_2t/\hbar}&\cdots\\
\vdots&\vdots&\ddots
\end{pmatrix}.
\end{aligned}
\end{equation}
The same statement is summarized without choosing a basis by
\begin{equation}
\boxed{U(t)=e^{-iHt/\hbar}},
\end{equation}
for a time-independent Hamiltonian.

\begin{classroomqa}
\lecturetimestamp{00:40:17}
\audiencequestion{Is this decomposition essentially a Fourier series, and is it analogous to decompositions used in heat conduction?}
\lecturerresponse{It is indeed a Fourier-like decomposition in time: each energy eigenvector carries a definite frequency. Fourier and Laplace methods also separate modes in many classical equations, including diffusion. The mathematical resemblance is real, but claims about which physical quantity corresponds to frequency must be checked case by case. In quantum mechanics the specific relation is fixed by the Schr\"odinger equation.}
\end{classroomqa}

\subsection{Combining Systems}

For two noninteracting systems \(A\) and \(B\), the combined Hamiltonian is the sum of the separate generators on the product space:
\begin{equation}
H_{AB}
=
H_A\otimes I_B+I_A\otimes H_B.
\end{equation}
If
\begin{equation}
H_A\ket{a}=E_a\ket{a},
\qquad
H_B\ket{b}=E_b\ket{b},
\end{equation}
then
\begin{equation}
H_{AB}\ket{a,b}=(E_a+E_b)\ket{a,b}.
\end{equation}
Energies, and therefore phase frequencies, add for separated noninteracting systems. An interaction contributes another term \(H_{\mathrm{int}}\), so this simple sum no longer gives the whole energy.

\begin{classroomqa}
\lecturetimestamp{00:56:20}
\audiencequestion{If two separate systems contain components with the same frequency, what happens when the systems are brought together?}
\lecturerresponse{For noninteracting systems, the combined energy eigenvalues are sums of the separate eigenvalues. Equality of two frequencies may produce degeneracy, but it does not by itself mix the corresponding eigenvectors. Mixing requires an interaction term or a choice of basis inside a degenerate subspace. Once the systems interact, the enlarged Hamiltonian must be diagonalized as a new problem.}
\end{classroomqa}

\section{Expectation Values and Commutators}

We have followed the state while keeping observables fixed. This is the Schr\"odinger picture. Let \(A\) be an observable with no explicit time dependence. Its expectation value is
\begin{equation}
\langle A\rangle_t
=
\bra{\psi(t)}A\ket{\psi(t)}.
\end{equation}
Differentiate both the bra and the ket:
\begin{equation}
\frac{d}{dt}\langle A\rangle_t
=
\bra{\dot\psi}A\ket{\psi}
+
\bra{\psi}A\ket{\dot\psi}.
\end{equation}
The Schr\"odinger equation and its Hermitian conjugate give
\begin{equation}
\ket{\dot\psi}
=-\frac{i}{\hbar}H\ket{\psi},
\qquad
\bra{\dot\psi}
=\frac{i}{\hbar}\bra{\psi}H.
\end{equation}
Hence
\begin{align}
\frac{d}{dt}\langle A\rangle_t
&=
\frac{i}{\hbar}\bra{\psi}HA\ket{\psi}
-\frac{i}{\hbar}\bra{\psi}AH\ket{\psi} \\
&=
\frac{i}{\hbar}\bra{\psi}(HA-AH)\ket{\psi}.
\end{align}
Define the commutator
\begin{equation}
[H,A]=HA-AH.
\end{equation}
Then the equation of motion for an expectation value is
\begin{equation}
\boxed{
\frac{d}{dt}\langle A\rangle_t
=
\frac{i}{\hbar}\langle[H,A]\rangle_t
}.
\end{equation}

\begin{figure}[htbp]
\centering
\includegraphics[width=0.74\textwidth]{lecture_09_figure_07.png}
\caption{The expectation-value equation and the commutator \(AB-BA=[A,B]\).}
\lectureframe{00:51:05}
\end{figure}

\begin{classroomqa}
\lecturetimestamp{00:49:13}
\audiencequestion{Why is the commutator in this equation itself inside an expectation value?}
\lecturerresponse{The commutator is an operator. Any operator placed between \(\bra{\psi}\) and \(\ket{\psi}\) has an expectation value, so \(\bra{\psi}[H,A]\ket{\psi}\) is simply the average of that commutator in the current state.}
\end{classroomqa}

\subsection{Conservation Laws}

Set \(A=H\). An operator always commutes with itself:
\begin{equation}
[H,H]=HH-HH=0.
\end{equation}
Therefore
\begin{equation}
\boxed{
\frac{d}{dt}\langle H\rangle_t=0
}.
\end{equation}
The expectation value of the Hamiltonian is conserved. This is the promised structural reason for identifying \(H\) as energy. Every isolated system has a Hamiltonian, and every time-independent Hamiltonian supplies this conservation law.

More generally,
\begin{equation}
[H,A]=0
\quad\Longrightarrow\quad
\frac{d}{dt}\langle A\rangle_t=0.
\end{equation}

\begin{classroomqa}
\lecturetimestamp{00:51:41}
\audiencequestion{What is the physical interpretation of the collection of all operators that commute with \(H\), and does it include quantities beyond energy?}
\lecturerresponse{They are the conserved observables of the system. Which ones occur depends on the Hamiltonian. A free particle conserves momentum, while a force can destroy that conservation. Energy is always present because \(H\) commutes with itself; additional commuting operators express additional symmetries and conservation laws.}
\end{classroomqa}

\section{One Spin in a Magnetic Field}

Consider a spin with Pauli operators
\begin{equation}
\boldsymbol{\sigma}
=(\sigma_1,\sigma_2,\sigma_3)
\end{equation}
in a fixed external magnetic field \(\mathbf B\). The magnetic moment is proportional to the spin. In the convention used here,
\begin{equation}
\boxed{
H=\frac{\mu}{2}\boldsymbol{\sigma}\cdot\mathbf B
}.
\end{equation}
The factor \(1/2\), the sign, and the charge-dependent scale may be absorbed into the definition of \(\mu\). The field is treated classically: a large stable magnet supplies \(\mathbf B\), and we do not model the quantum state of that magnet.

\begin{figure}[htbp]
\centering
\includegraphics[width=0.68\textwidth]{lecture_09_figure_08.png}
\caption{The magnetic-moment operator couples to the external magnetic field.}
\lectureframe{01:01:45}
\end{figure}

Choose the field along the third axis,
\begin{equation}
\mathbf B=B\hat{\mathbf z},
\qquad
H=\frac{\mu B}{2}\sigma_3.
\end{equation}
We shall track the expectation values of the three spin components. To shorten the equations, write
\begin{equation}
s_i(t)=\langle\sigma_i\rangle_t.
\end{equation}
The longitudinal component obeys
\begin{equation}
\dot s_3
=
\frac{i}{\hbar}\langle[H,\sigma_3]\rangle
=0,
\end{equation}
because \(H\) is proportional to \(\sigma_3\). The projection of the spin along the field is fixed. Geometrically, the tip of the average spin must remain on a cone of fixed opening angle about the field axis.

\begin{figure}[htbp]
\centering
\includegraphics[width=0.58\textwidth]{lecture_09_figure_09.png}
\caption{Fixed longitudinal spin leaves motion on a cone about the magnetic-field axis.}
\lectureframe{01:04:45}
\end{figure}

For the first transverse component,
\begin{equation}
\dot s_1
=
\frac{i\mu B}{2\hbar}
\langle[\sigma_3,\sigma_1]\rangle.
\end{equation}
The commutator can be checked directly from
\begin{equation}
\begin{aligned}
\sigma_3&=
\begin{pmatrix}1&0\\0&-1\end{pmatrix},
&
\sigma_1&=
\begin{pmatrix}0&1\\1&0\end{pmatrix},\\[0.5em]
\sigma_2&=
\begin{pmatrix}0&-i\\i&0\end{pmatrix}.
\end{aligned}
\end{equation}
Indeed,
\begin{equation}
\sigma_3\sigma_1=i\sigma_2,
\qquad
\sigma_1\sigma_3=-i\sigma_2,
\end{equation}
so
\begin{equation}
[\sigma_3,\sigma_1]=2i\sigma_2.
\end{equation}
Consequently,
\begin{equation}
\dot s_1=-\frac{\mu B}{\hbar}s_2.
\end{equation}
Likewise, using \([\sigma_3,\sigma_2]=-2i\sigma_1\),
\begin{equation}
\dot s_2=\frac{\mu B}{\hbar}s_1.
\end{equation}

\begin{figure}[htbp]
\centering
\includegraphics[width=0.72\textwidth]{lecture_09_figure_10.png}
\caption{The Pauli commutator calculation leading to the transverse equation of motion.}
\lectureframe{01:09:35}
\end{figure}

Let
\begin{equation}
\Omega=\frac{\mu B}{\hbar}.
\end{equation}
The three equations are
\begin{equation}
\dot s_1=-\Omega s_2,
\qquad
\dot s_2=\Omega s_1,
\qquad
\dot s_3=0.
\end{equation}
A solution with transverse amplitude \(s_\perp\) is
\begin{equation}
s_1(t)=s_\perp\cos(\Omega t+\varphi),
\qquad
s_2(t)=s_\perp\sin(\Omega t+\varphi),
\end{equation}
with \(s_3\) constant.

\begin{figure}[htbp]
\centering
\includegraphics[width=0.64\textwidth]{lecture_09_figure_11.png}
\caption{The transverse spin components form a sine--cosine pair with angular frequency \(\mu B/\hbar\).}
\lectureframe{01:11:45}
\end{figure}

The average spin therefore precesses around the magnetic field like a gyroscope. Choosing the field along \(z\) was only a convenience. For a field in any direction, the component parallel to \(\mathbf B\) is conserved and the perpendicular component rotates around it.

\begin{classroomqa}
\lecturetimestamp{01:13:02}
\audiencequestion{Should the sine and cosine solution contain an additional phase?}
\lecturerresponse{Yes. The phase \(\varphi\) specifies the initial direction in the transverse plane, or equivalently the choice of the instant called \(t=0\). Setting it to zero merely chooses the spin to begin along one particular transverse axis.}
\end{classroomqa}

\subsection{Preparation and Measurement}

Suppose an ensemble is prepared in the \(+1\) eigenstate of \(\sigma_1\). At \(t=0\), a measurement of \(\sigma_1\) is certain to return \(+1\), while measurements of \(\sigma_2\) or \(\sigma_3\) give equal probabilities for \(\pm1\). After a quarter precession period, the state has rotated to the \(+1\) eigenstate of \(\sigma_2\): \(\sigma_2\) is then certain and \(\sigma_1\) is maximally uncertain. After another quarter period it reaches the \(-1\) eigenstate of \(\sigma_1\). The rotating expectation values therefore have direct operational meaning.

\begin{classroomqa}
\lecturetimestamp{01:14:04}
\audiencequestion{If expectation values require many electrons, will the members of the ensemble not carry different phases and wash out the precession?}
\lecturerresponse{Not when the same experiment is performed on every member. Identical preparation gives every electron the same initial state and hence the same phase evolution. If their initial phases are uncontrolled, the ensemble is not a pure preparation; it is a mixed state, and dephasing can reduce the transverse average.}
\end{classroomqa}

\section{Two Coupled Spins}

Now consider two spins. Merely saying that they are entangled is not enough to predict what they will do. Dynamics requires a Hamiltonian.

\begin{classroomqa}
\lecturetimestamp{01:16:49}
\audiencequestion{If two electrons are entangled so that opposite results occur along any common axis, does that mean their motions are always \(180^\circ\) out of phase?}
\lecturerresponse{No dynamical conclusion follows from entanglement alone. The stated anticorrelation describes the singlet state, but its subsequent evolution depends on the energy operator. We must first specify whether there is an external field, an interaction between the spins, or both.}
\end{classroomqa}

Take the simple isotropic interaction in which each spin supplies the effective field felt by the other:
\begin{equation}
\boxed{
H=J\,\boldsymbol{\sigma}\cdot\boldsymbol{\tau}
=
J\left(
\sigma_1\tau_1+\sigma_2\tau_2+\sigma_3\tau_3
\right)
}.
\end{equation}
The operators \(\boldsymbol{\sigma}\) act on the first spin and \(\boldsymbol{\tau}\) on the second. The coupling \(J\) sets the energy scale; its value and sign are left unspecified.

\begin{figure}[htbp]
\centering
\includegraphics[width=0.75\textwidth]{lecture_09_figure_12.png}
\caption{The isotropic two-spin interaction \(H\propto\boldsymbol{\sigma}\cdot\boldsymbol{\tau}\).}
\lectureframe{01:19:20}
\end{figure}

\begin{classroomqa}
\lecturetimestamp{01:17:49}
\audiencequestion{Would placing two entangled electrons in a uniform magnetic field necessarily disentangle them?}
\lecturerresponse{Not necessarily. A field changes the Hamiltonian and may tend to align the spins, especially if it is strong, but unitary evolution in a uniform field does not by itself guarantee loss of entanglement. The answer depends on the complete Hamiltonian and the initial state.}
\end{classroomqa}

\subsection{Triplet and Singlet Energies}

The two-spin Hilbert space has dimension four. The interaction is diagonal in the triplet--singlet basis:
\begin{align}
\ket{T_+}&=\ket{\uparrow\uparrow}, \\
\ket{T_0}&=
\frac{\ket{\uparrow\downarrow}+\ket{\downarrow\uparrow}}{\sqrt2}, \\
\ket{T_-}&=\ket{\downarrow\downarrow}, \\
\ket{S}&=
\frac{\ket{\uparrow\downarrow}-\ket{\downarrow\uparrow}}{\sqrt2}.
\end{align}
The three triplet states share one energy \(E_T\); the singlet has another energy \(E_S\). Thus there are four independent eigenvectors but only two distinct energy levels. The triplet level is threefold degenerate.

\begin{figure}[htbp]
\centering
\includegraphics[width=0.74\textwidth]{lecture_09_figure_13.png}
\caption{The coupled-spin Hamiltonian has a degenerate triplet level and a separate singlet level.}
\lectureframe{01:20:30}
\end{figure}

\begin{classroomqa}
\lecturetimestamp{01:21:16}
\audiencequestion{What happens if the initial state is already the singlet energy eigenstate?}
\lecturerresponse{It remains the singlet and acquires only the phase \(e^{-iE_St/\hbar}\). As with any single energy eigenstate, its measurable properties do not change merely because the whole ket is multiplied by that common phase.}
\end{classroomqa}

\begin{classroomqa}
\lecturetimestamp{01:21:54}
\audiencequestion{Does an external magnetic field supply energy to the two-spin system during this calculation?}
\lecturerresponse{There is no external field in the example being solved. The only field under discussion is the mutual interaction represented by \(J\,\boldsymbol{\sigma}\cdot\boldsymbol{\tau}\). Adding an external field would add another term to \(H\) and define a different problem.}
\end{classroomqa}

\subsection{Begin with One Spin Up and One Down}

An energy eigenstate would only gather a phase, so choose a product state that is not an energy eigenstate:
\begin{equation}
\ket{\psi(0)}=\ket{\uparrow\downarrow}.
\end{equation}
Decompose it into the two relevant eigenvectors:
\begin{equation}
\boxed{
\ket{\uparrow\downarrow}
=
\frac{1}{\sqrt2}\ket{T_0}
+
\frac{1}{\sqrt2}\ket{S}
}.
\end{equation}

\begin{figure}[htbp]
\centering
\includegraphics[width=0.76\textwidth]{lecture_09_figure_14.png}
\caption{The product state \(\ket{\uparrow\downarrow}\) decomposed into triplet and singlet components.}
\lectureframe{01:24:30}
\end{figure}

Each energy component now receives its own phase:
\begin{equation}
\ket{\psi(t)}
=
\frac{e^{-iE_Tt/\hbar}}{\sqrt2}\ket{T_0}
+
\frac{e^{-iE_St/\hbar}}{\sqrt2}\ket{S}.
\end{equation}

\begin{classroomqa}
\lecturetimestamp{01:25:59}
\audiencequestion{Which is larger, the triplet energy \(E_T\) or the singlet energy \(E_S\)?}
\lecturerresponse{That depends on the sign of the coupling. For \(H=J\,\boldsymbol{\sigma}\cdot\boldsymbol{\tau}\) with \(J>0\), the singlet is lower and the triplet is higher; reversing the sign reverses the ordering. Adding a constant to both energies changes only an overall phase, so one may set either level to zero without changing the transition probabilities.}
\end{classroomqa}

Return to the product basis:
\begin{align}
\ket{\psi(t)}
&=
\frac12
\left(
e^{-iE_Tt/\hbar}+e^{-iE_St/\hbar}
\right)\ket{\uparrow\downarrow} \\
&\quad+
\frac12
\left(
e^{-iE_Tt/\hbar}-e^{-iE_St/\hbar}
\right)\ket{\downarrow\uparrow}.
\end{align}

\begin{figure}[htbp]
\centering
\includegraphics[width=0.76\textwidth]{lecture_09_figure_15.png}
\caption{The evolved state reassembled in the \(\ket{\uparrow\downarrow}\), \(\ket{\downarrow\uparrow}\) basis.}
\lectureframe{01:30:40}
\end{figure}

At \(t=0\), the coefficient of \(\ket{\downarrow\uparrow}\) vanishes, as it must. At later times the relative phase between the singlet and triplet components makes this coefficient grow. Define the mean energy and splitting,
\begin{equation}
\overline E=\frac{E_T+E_S}{2},
\qquad
\Delta E=E_T-E_S.
\end{equation}
Factoring out the mean-energy phase gives
\begin{equation}
\boxed{
\begin{aligned}
\ket{\psi(t)}
&=e^{-i\overline E t/\hbar}\Biggl[
\cos\!\left(\frac{\Delta E\,t}{2\hbar}\right)
\ket{\uparrow\downarrow}
\\
&\qquad
-i\sin\!\left(\frac{\Delta E\,t}{2\hbar}\right)
\ket{\downarrow\uparrow}
\Biggr]
\end{aligned}
}.
\end{equation}

\begin{figure}[H]
\centering
\includegraphics[width=0.76\textwidth]{lecture_09_figure_16.png}
\caption{The mean energy becomes a common phase, while the singlet--triplet splitting appears in the sine and cosine amplitudes.}
\lectureframe{01:35:15}
\end{figure}

The factor \(e^{-i\overline E t/\hbar}\) is common to the entire ket and drops out of every probability. The measurable probabilities are
\begin{equation}
P_{\uparrow\downarrow}(t)
=
\cos^2\!\left(\frac{\Delta E\,t}{2\hbar}\right),
\qquad
P_{\downarrow\uparrow}(t)
=
\sin^2\!\left(\frac{\Delta E\,t}{2\hbar}\right).
\end{equation}
The state oscillates from \(\ket{\uparrow\downarrow}\) to \(\ket{\downarrow\uparrow}\) and back, at a rate determined only by the energy difference. At intermediate times both amplitudes are nonzero, and except at special instants the two spins are entangled.

The calculation illustrates a general method. Choose the basis in which the physical question is posed, expand the initial state in eigenvectors of the Hamiltonian, let each eigenvector acquire its own phase, and then recombine the pieces in the measurement basis. Quantum dynamics becomes systematic once the Hamiltonian has been diagonalized; all observable motion is carried by relative phase.
